% BEGIN R28_REGIONAL_CLOSURES
% Insert after REG01. This is not a standalone document.
% Provenance: U012_biografia_estructural_pi.tex, lines 175--324.
% Oriented witnesses: verificar_cierres_regionales.py; same problem,
% representatives different from the walks printed in U012.
% Emitter SHA-256: 42fdacab582de8e2d2cb7dfafc6204d200aef01249b3dcc37949e39507f066d8.
\subsubsection{Block multigraph and five minimal regional closures}
\label{regcl27:seccion}

The \(468\) emissions in \(\mathcal U_6\), produced by the
APP--TRIT--TPK emitter, determine a block multigraph. Identifying its
regional anchors makes it possible to solve a joint edge-covering
problem: to traverse the APP, propagation, autoscale, and each closure
region's incidences within a single closed walk. The local six-junction
cycle of Section~\ref{sec:rec27-empalmes} and these covers retain their
respective traversal conditions.

\paragraph{Block graph and exact phase quotient.}
For \(U=(U_1,\ldots,U_6)\in\mathcal U_6\), define
\begin{equation}
 A(U)=(U_1,U_2,U_3),\qquad B(U)=(U_4,U_5,U_6).
 \label{regcl27:mitades}
\end{equation}
The undirected multigraph \(\widetilde G\) has as vertices the triples
occurring as \(A(U)\) or \(B(U)\), and one edge labelled by each \(U\),
incident at those two triples. All \(468\) labels are retained:
distinct emissions are not identified merely because they share
endpoints or because one interchanges the halves of another.

Internal phase acts on each triple by cyclic rotation:
\begin{equation}
 (a,b,c)\sim(b,c,a)\sim(c,a,b),\qquad
 \operatorname{can}(a,b,c)
 =\min_{\rm lex}\{(a,b,c),(b,c,a),(c,a,b)\}.
 \label{regcl27:fase}
\end{equation}
The graph \(G\) replaces the endpoints of edge \(U\) by
\(\operatorname{can}A(U)\) and \(\operatorname{can}B(U)\), retaining
its label. The equivalence acts precisely by the three cyclic rotations
of each vertex; every emission retains its identity as an edge.

Enumeration from the same emitter and component searches give
\begin{equation}
\begin{array}{c|r|l}
 & |V|&\text{component sizes}\\ \hline
 \widetilde G&96&(28,28,28,4,4,4)\\
 G&32&(28,4)
\end{array}
\label{regcl27:componentes}
\end{equation}
These data are obtained by collecting distinct endpoints, adding each
incidence in both directions, and traversing each unvisited component
by breadth-first search. Adjacency may be deduplicated when counting
components, but parallel labels are retained when searching for
closures with prescribed edges.

\paragraph{APP anchor and regional numbering.}
The three common edges are
\begin{equation}
\begin{aligned}
 U_{\rm APP}&=903|850|850|252|109|252,\\
 U_e&=850|963|890|149|252|212,\\
 U_\varphi&=830|943|830|109|252|252.
\end{aligned}
\label{regcl27:anclajes}
\end{equation}
Fix the endpoint
\(v_{\rm APP}=\operatorname{can}A(U_{\rm APP})=850|850|903\).
For the five already constructed regions, use the following order,
which places the transverse region in the fifth row:
\begin{equation}
\begin{array}{c|l|c|r}
 i&U_\pi^{(i)}&\rho_i&m_i\\ \hline
 1&810|923|870|169|272|272&++&144\\
 2&810|983|810|169|272|272&++&144\\
 3&870|923|810|169|272|272&++&144\\
 4&870|923|810|418|674|521&--&144\\
 5&501|614|498|169|272|272&\perp&432
\end{array}
\label{regcl27:regiones}
\end{equation}
This numbering distinguishes the three \(++\) regions, the mutated
return \( -- \), and the transverse region \(\perp\). Accordingly,
\(\rho=(++ ,++ ,++ ,--,\perp)\) and
\(m=(144,144,144,144,432)\) in this order; their sum remains \(1008\).

\paragraph{Five complete witnesses with traversal directions.}
In each table, row \(j\) means traversing the edge labelled \(e_j\)
from \(v_{j-1}\) to \(v_j\). Endpoints are canonical phase representatives;
since \(G\) is undirected, the pair \((v_{j-1},v_j)\) specifies
the traversal direction of the labelled edge.
Each table contains exactly eight edges and satisfies
\(v_0=v_8=v_{\rm APP}\).

The breadth-first search described in the proof yields the following
witnesses, one for each set of four target edges. Minimality concerns
length; several walks may attain that length.
\paragraph{Witness \(i=1\).}
\begin{center}
\small
\begin{tabular}{rlll}
\hline
\(j\)&\(v_{j-1}\)&\(e_j\)&\(v_j\)\\
\hline
1&\(850|850|903\)&\(109|252|252|850|903|850\)&\(109|252|252\)\\
2&\(109|252|252\)&\(109|252|252|810|923|870\)&\(810|923|870\)\\
3&\(810|923|870\)&\(810|923|870|169|272|272\)&\(169|272|272\)\\
4&\(169|272|272\)&\(169|272|272|850|963|890\)&\(850|963|890\)\\
5&\(850|963|890\)&\(850|963|890|149|252|212\)&\(149|252|212\)\\
6&\(149|252|212\)&\(149|252|212|830|943|830\)&\(830|830|943\)\\
7&\(830|830|943\)&\(830|943|830|109|252|252\)&\(109|252|252\)\\
8&\(109|252|252\)&\(903|850|850|252|109|252\)&\(850|850|903\)\\
\hline
\end{tabular}
\end{center}

\paragraph{Witness \(i=2\).}
\begin{center}
\small
\begin{tabular}{rlll}
\hline
\(j\)&\(v_{j-1}\)&\(e_j\)&\(v_j\)\\
\hline
1&\(850|850|903\)&\(169|272|272|850|903|850\)&\(169|272|272\)\\
2&\(169|272|272\)&\(169|272|272|810|983|810\)&\(810|810|983\)\\
3&\(810|810|983\)&\(810|983|810|169|272|272\)&\(169|272|272\)\\
4&\(169|272|272\)&\(169|272|272|850|963|890\)&\(850|963|890\)\\
5&\(850|963|890\)&\(850|963|890|149|252|212\)&\(149|252|212\)\\
6&\(149|252|212\)&\(149|252|212|830|943|830\)&\(830|830|943\)\\
7&\(830|830|943\)&\(830|943|830|109|252|252\)&\(109|252|252\)\\
8&\(109|252|252\)&\(903|850|850|252|109|252\)&\(850|850|903\)\\
\hline
\end{tabular}
\end{center}

\paragraph{Witness \(i=3\).}
\begin{center}
\small
\begin{tabular}{rlll}
\hline
\(j\)&\(v_{j-1}\)&\(e_j\)&\(v_j\)\\
\hline
1&\(850|850|903\)&\(169|272|272|850|903|850\)&\(169|272|272\)\\
2&\(169|272|272\)&\(169|272|272|850|963|890\)&\(850|963|890\)\\
3&\(850|963|890\)&\(850|963|890|149|252|212\)&\(149|252|212\)\\
4&\(149|252|212\)&\(149|252|212|870|923|810\)&\(810|870|923\)\\
5&\(810|870|923\)&\(870|923|810|169|272|272\)&\(169|272|272\)\\
6&\(169|272|272\)&\(169|272|272|830|943|830\)&\(830|830|943\)\\
7&\(830|830|943\)&\(830|943|830|109|252|252\)&\(109|252|252\)\\
8&\(109|252|252\)&\(903|850|850|252|109|252\)&\(850|850|903\)\\
\hline
\end{tabular}
\end{center}

\paragraph{Witness \(i=4\).}
\begin{center}
\small
\begin{tabular}{rlll}
\hline
\(j\)&\(v_{j-1}\)&\(e_j\)&\(v_j\)\\
\hline
1&\(850|850|903\)&\(169|272|272|850|903|850\)&\(169|272|272\)\\
2&\(169|272|272\)&\(169|272|272|850|963|890\)&\(850|963|890\)\\
3&\(850|963|890\)&\(850|963|890|149|252|212\)&\(149|252|212\)\\
4&\(149|252|212\)&\(149|252|212|870|923|810\)&\(810|870|923\)\\
5&\(810|870|923\)&\(870|923|810|418|674|521\)&\(418|674|521\)\\
6&\(418|674|521\)&\(418|674|521|830|943|830\)&\(830|830|943\)\\
7&\(830|830|943\)&\(830|943|830|109|252|252\)&\(109|252|252\)\\
8&\(109|252|252\)&\(903|850|850|252|109|252\)&\(850|850|903\)\\
\hline
\end{tabular}
\end{center}

\paragraph{Witness \(i=5\).}
\begin{center}
\small
\begin{tabular}{rlll}
\hline
\(j\)&\(v_{j-1}\)&\(e_j\)&\(v_j\)\\
\hline
1&\(850|850|903\)&\(109|252|252|850|903|850\)&\(109|252|252\)\\
2&\(109|252|252\)&\(109|252|252|501|614|498\)&\(498|501|614\)\\
3&\(498|501|614\)&\(501|614|498|169|272|272\)&\(169|272|272\)\\
4&\(169|272|272\)&\(169|272|272|850|963|890\)&\(850|963|890\)\\
5&\(850|963|890\)&\(850|963|890|149|252|212\)&\(149|252|212\)\\
6&\(149|252|212\)&\(149|252|212|830|943|830\)&\(830|830|943\)\\
7&\(830|830|943\)&\(830|943|830|109|252|252\)&\(109|252|252\)\\
8&\(109|252|252\)&\(903|850|850|252|109|252\)&\(850|850|903\)\\
\hline
\end{tabular}
\end{center}

\begin{theorem}[Minimality of the five regional edge covers]
\label{regcl27:minimalidad}
For each \(i\in\{1,\ldots,5\}\), among closed walks in \(G\)
traversing every edge of
\[
 E_i^*=\{U_{\rm APP},U_e,U_\varphi,U_\pi^{(i)}\}
\]
at least once, the minimum length, counting traversals with multiplicity,
is eight. Vertices and edges may be repeated.
\end{theorem}

\begin{proof}
Every closed walk traversing \(U_{\rm APP}\) passes through
\(v_{\rm APP}\). Cyclic rotation of its step list allows it to start at
that vertex without changing length or coverage. It therefore suffices
to solve the problem with a fixed starting vertex.

The search space is the finite set
\[
 \mathscr S_i=V(G)\times\mathcal P(E_i^*),
 \qquad |\mathscr S_i|\le32\cdot16=512.
\]
The initial state is \(s_0=(v_{\rm APP},\varnothing)\), and the terminal
state is \(s_*=(v_{\rm APP},E_i^*)\). Traversing an edge \(e\)
incident at \(v\) towards its other endpoint \(v'\) updates the state by
\begin{equation}
 (v,M)\longmapsto\bigl(v',M\cup(E_i^*\cap\{e\})\bigr).
 \label{regcl27:transicion-bfs}
\end{equation}
The mask records edge labels, not just endpoints; distinct parallel
edges must therefore not be collapsed.

Every transition has unit cost. A FIFO queue initialized at \(s_0\)
extracts states in increasing distance, marks each state on its first
visit, and stores its predecessor and traversed label.
Equivalently, with \(\mathcal F_0=\{s_0\}\), the new layers are
\[
 \mathcal F_{k+1}
 =\operatorname{Succ}(\mathcal F_k)
   \setminus\bigcup_{j=0}^{k}\mathcal F_j.
\]
Induction shows that \(\mathcal F_k\) contains exactly the states whose
minimum distance from \(s_0\) is \(k\). Finite enumeration over the
specified emitter gives, for each of the five regions,
\[
 s_*\notin\bigcup_{k=0}^{7}\mathcal F_k,\qquad s_*\in\mathcal F_8,
 \qquad (\ell_1,\ell_2,\ell_3,\ell_4,\ell_5)=(8,8,8,8,8).
\]
The tables exhibit the five upper bounds, including all four required
labels in each case. Exhaustive exploration of the preceding layers
proves the lower bound. Together they establish minimality for the
declared problem.
\end{proof}

Length eight characterizes the joint coverage of four incidences in
each closure region. The quotient retains emission labels and their
junctions between phase classes; enriched temporal evolution also
preserves orientation and memory. This description links all five
regions to the same propagation, autoscale, and APP anchors and
complements the local six-junction cycle. Coinductive prolongation
retains its operators and cylinders together with these orientation
and memory data.
% END R28_REGIONAL_CLOSURES
